\documentclass[12pt, a4paper]{article}

\usepackage[utf8]{inputenc}
\usepackage[T1]{fontenc}
\usepackage[italian]{babel}
\usepackage{geometry}

\usepackage[table]{xcolor} 
\usepackage{booktabs}      
\usepackage{tabularx} 
\usepackage{enumitem} 

\geometry{a4paper, margin=1.5cm} 

\definecolor{headerblue}{RGB}{31, 78, 121}
\definecolor{rowgray}{RGB}{245, 245, 245}

\begin{document}

\section*{Project Classification}

\vspace{1em}

\begin{table}[htbp]
\centering
\small
\renewcommand{\arraystretch}{1.6} 
\rowcolors{2}{white}{rowgray}     

\begin{tabularx}{\textwidth}{l c c c X c} 
\toprule
\rowcolor{headerblue} 
\textcolor{white}{\textbf{CLASSE}} & 
\textcolor{white}{\textbf{DURATA}} & 
\textcolor{white}{\textbf{RISCHIO}} & 
\textcolor{white}{\textbf{COMPLESSITÀ}} & 
\textcolor{white}{\textbf{TECNOLOGIA}} & 
\textcolor{white}{\textbf{PROBABILITÀ}} \\
\midrule

\textbf{Type A} & >24 months   & ALTA        & ALTA        & Breakthrough / Experimental & CERTAMENTE \\
\textbf{Type B} & 12-24 months & MEDIO       & MEDIO       & Current / Leading Edge      & MOLTO PROBABILE \\
\textbf{Type C} & 6-12 months  & BASSO       & BASSO       & Best of Breed / Proven      & PROBABILE \\
\textbf{Type D} & 3-6 months   & MOLTO BASSO & MOLTO BASSO & Standard / Off-the-Shelf    & POCO PROBABILE \\
\textbf{Type E} & <3 months    & MINIMO      & MINIMO      & Practical / Legacy          & NULLO \\

\bottomrule
\end{tabularx}
\end{table}

\vspace{2em} 

\subsection*{Glossario delle Tecnologie}
\begin{itemize}[label=\textbullet, itemsep=0.5em]
    \item \textbf{Breakthrough / Experimental:} Tecnologie all'avanguardia o in fase di test sperimentale. Offrono enormi vantaggi competitivi ma comportano un rischio di stabilità altissimo.
    \item \textbf{Current / Leading Edge:} Tecnologie moderne e ampiamente adottate dal mercato attuale. Rappresentano il compromesso ideale tra innovazione e affidabilità strutturale.
    \item \textbf{Best of Breed / Proven:} I leader storici e super-collaudati per una specifica funzione. Soluzioni mature, stabili e con ampio supporto garantito.
    \item \textbf{Standard / Off-the-Shelf:} Prodotti software commerciali o open-source "pronti all'uso" e pacchettizzati. Richiedono configurazione anziché sviluppo da zero.
    \item \textbf{Practical / Legacy:} Tecnologie basilari o sistemi consolidati da anni. Livello di innovazione assente, ma comportamento prevedibile al 100\% e rischio nullo.
\end{itemize}

\end{document}